% =====================================================================
% WORKED FIGURE CAPTIONS AND METHOD NOTES — staging file.
%
% Why this file exists: figures are not yet placed in HIKET_main_manuscript.tex
% (deliberately held until the storyline settles), but some captions have already
% been worked out and are too developed to re-derive. They live here until the
% figures are placed, then move into the manuscript and this file goes away.
%
% Rescued 2026-09-03 from HIKET_storyline_note.tex (retired 2026-09-02) — that
% file was the ONLY home of the F14 method note.
%
% ⚠ NUMBERS HAVE BEEN STRIPPED, NOT UPDATED. Every figure quoted in the original
% was measured on a superseded run, and three of its claims have since died:
% the "simple models show no such ridge" (falsified by S13 — with xi included the
% trio rides a STRONGER ridge), the ll costs 1.6/14.7 at the published MRT, and
% Yasso20's published 19.0 yr (a chimera: Yasso20 fractions paired with Yasso15
% rates; corrected to 21.85 point / 22.13 posterior). Refill from the current run
% before use. The prose below is structural and survives; the digits did not.
% =====================================================================

% ---------------------------------------------------------------------
% F14 — what a longer residence time would cost
% ---------------------------------------------------------------------
\begin{figure}[htbp]\centering
\includegraphics[width=0.98\linewidth]{\fig/F14_mrt_ridge.png}
\caption{\textbf{What a longer residence time would cost.} \emph{Top:} best attainable fit over
(intrinsic MRT, $\sinput$) for each Yasso model. \textbf{Colours are the likelihood alone --- the
fit to the data, priors excluded --- blue is the best fit, through green and yellow to red at the
worst. The white contour lines are the posterior --- likelihood \emph{and} priors together.} Colour
is the log-likelihood deficit from the best fit found, clamped at 25 units. Dashed line: published
MRT; white dot: posterior median. Regions the sampler visited fewer than five times are left blank.
\emph{Bottom:} first-order variance indices of MRT. The colour peak lies up and to the left of the
contours in all three panels: the data alone want an \emph{even shorter} residence time, and it is
the prior that pulls the answer back toward the published value.}\label{fig:ridge}
\end{figure}

\begin{methodnote}
\textbf{Reading the two layers.} The panels superimpose two different things on the same axes, and
keeping them apart is the whole point of the figure.

\emph{The white contour lines are the posterior} --- where the calibration actually ended up. That
location is set by three ingredients at once: how well the data are fitted, what the priors say, and
how many parameter combinations happen to land there (many parameters map onto the same residence
time, and that alone attracts probability mass).

\emph{The colours are the likelihood only} --- the best fit to the data achievable in that cell,
with the priors switched off. They answer a different question: not ``where did we end up'', but
``how well \emph{could} the data be fitted here''.

So the offset between the bluest colour and the innermost contour \textbf{measures how far the prior
pulls the answer away from where the data alone would put it}. Here that offset always points the
same way: the likelihood wants shorter residence times and larger inputs than we report.

\emph{Three honest limits.} (i) The colours are prior-free in \emph{value} but not in
\emph{coverage}: each cell shows the best of the draws the sampler happened to visit, and the
sampler was steered by the priors. It is therefore a lower bound on the true best fit, and it is
weakest exactly where draws are thin --- and thinnest, for Yasso07, to the point of being absent
around the published value. A blank cell means \emph{not visited}, never \emph{cannot be fitted},
and Yasso07's blank is precisely why a profile likelihood is needed. (ii) The colour surface is
estimated on a coarse grid --- one whose cells hold enough draws to support a maximum --- and then
bilinearly interpolated for display, so it looks continuous without claiming resolution the draws
cannot support. (iii) Colour is the log-likelihood \emph{deficit} from the best fit found, clamped
at 25 units: a deficit past that is decisively rejected, so resolving 30 from 200 would convey
nothing and would flatten the few units of structure along the ridge that actually matter.
\end{methodnote}

\begin{methodnote}
\textbf{Which published value.} Costs are measured against the published \emph{point}
parameterisation. Yasso15 and Yasso20 also publish posterior \emph{samples}, and because residence
time is a nonlinear many-to-one function of the parameters, the median residence time of that sample
is not the residence time of the median parameters. \textbf{The two comparators must never be
mixed}; the figure uses the point throughout.
\end{methodnote}

\begin{provisional}
\textbf{To be redone before submission.} Replace the sampled colour map with a true \emph{profile
likelihood} --- an optimisation over the remaining parameters at each grid node --- which removes
the prior from the coverage as well as from the value, prices the region beyond the posterior's
reach, and is smooth enough to interpolate. Roughly 7\,h per model on a $20\times20$ grid;
\textbf{Yasso07 first}, since it is the model for which no sampled estimate exists at all.
\end{provisional}
